% SOURCE_AUTHORITY: sga5_fr_workpass.tex, lines 12133--12168 (LF-normalized unit).
% SOURCE_SHA256: 791F4EFFC5E02832D5D77ED03518C8156D6F07E4C8238B03545DB93D883FBB28
\subsection*{9. Propriedades funtoriais}

a) Sejam $A$ e $B$ dois anéis (não necessariamente comutativos), e seja $u:A\to B$ um homomorfismo. Todo $B$-módulo à esquerda (resp. à direita) $M$ pode ser dotado de uma estrutura de $A$-módulo à esquerda (resp. à direita), que se denotará por $M_{[u]}$. Em particular, $B_{[u]}$ é um $A$-módulo à direita.

O funtor
\[
E\longmapsto B_{[u]}\otimes_A E
\]
permite definir o funtor exato
\[
D(A)_{\parf}\longrightarrow D(B)_{\parf}
\]
e, consequentemente, o homomorfismo:
\[
K^{\bullet}(A)\xrightarrow{u^*}K^{\bullet}(B).
\]
Obtém-se, de fato, um funtor covariante da categoria dos anéis (resp. dos anéis comutativos) para a categoria dos grupos abelianos (resp. dos anéis comutativos).

b) Suponhamos agora que o $A$-módulo à esquerda $B_{[u]}$ seja pseudocoerente. Resulta, portanto, que o funtor
\[
F\longmapsto F_{[u]}
\]
transforma um $B$-módulo livre finitamente gerado em um $A$-módulo pseudocoerente. Utilizando a proposição 5.1, obtém-se um funtor exato:
\[
D^b(B)_{\coh}\longrightarrow D^b(A)_{\coh}
\]
e, consequentemente, um homomorfismo canônico
\[
K_\bullet(B)\xrightarrow{u_*}K_\bullet(A).
\]
Os homomorfismos $u_*$, $u^*$ verificam a seguinte identidade (a fórmula de projeção):
\[
u_*(u^*(\alpha)\cdot \beta)=\alpha\,(u_*(\beta))
\quad\text{para }\alpha\in K^{\bullet}(A),\ \beta\in K_\bullet(B).
\]
Esta fórmula simples será demonstrada de forma mais geral nas exposições do seminário sobre Riemann--Roch (SGA 6 IV 2.11.1.2).
